The proposed method, and its submodules, were extensively evaluated as presented below. 

\subsection{Training Methodology}\label{sec:training_dataset}
A composite dataset of both real and simulated images was collected for training the VAE network. Real data was collected with an Intel RealSense D$455$ depth camera. This dataset consists of images in confined spaces, indoor rooms, long corridors, and outdoor environments with trees. Thin obstacles such as tree branches, rods, and poles are manually labeled in this dataset. We also utilize images from the NYU Depth Dataset v2~\cite{silberman2012nyu_depth}. Simulation images are collected using both Gazebo Classic~\cite{gazbeoclassic} and Isaac Gym Simulators~\cite{isaac_gym}. Isaac Gym offers a segmentation camera providing instance segmentation masks for the corresponding depth images. Simulated meshes having cross sections below a size of $5~\textrm{cm}$ are assigned semantic instance IDs. The segmentation camera allows rendering images with pixel values equal to the ID of the semantic occupying that pixel. The aggregated dataset $\mathbf{X}$ consists of $\sim 66,000$ images, of which $\sim 35,000$ are simulated. $21,000$ simulated images have semantic labels, and $\sim 1,000$ real images are labeled. The training and validation sets are split with a $80\%:20\%$ ratio from this dataset. We train the network using the Adam Optimizer~\cite{kingma2015adam} with a learning rate of $10^{-4}$ for $40~\textrm{epochs}$.

%%%%%%%%%%%%%%%%%%%%%%%%%%%%%%%%%%%%%%%%%%%%%%%%%%%%%%%%%%%%%%%%%
\begin{figure}[h]
\includegraphics[width=\columnwidth]{figures/DataCollectionDiagram_v5.png}
\vspace{-4ex}
\caption{Instances of depth images (1b, 2b) collected for the \ac{sevae} training from the real world (1a) and the IsaacGym (2a) simulator are shown. RGB images (1a) from a collected dataset (if available) are used to create semantic labels (1c) for the aligned depth images. The objects beyond the sensing range of the depth camera are not labeled. The IsaacGym simulator provides semantic labels for thin obstacles (2c). Simulation data for training the \ac{cpn} is collected in Gazebo Classic (3) with a simulated drone in an obstacle-filled environment. Randomized action sequences are executed to collect a collision dataset.}
\vspace{-1ex}
\label{fig:data_collection_diagram}
\end{figure}
%%%%%%%%%%%%%%%%%%%%%%%%%%%%%%%%%%%%%%%%%%%%%%%%%%%%%%%%%%%%%%%%%

Collision datasets for training the \ac{cpn} are collected using the Gazebo-based RotorS simulator~\cite{Furrer2016rotors} with a simulated robot model. An environment consisting of obstacles having different shapes and sizes is constructed and thin obstacles with cross sections smaller than $5~\textrm{cm}$ are introduced as shown in Figure~\ref{fig:data_collection_diagram}. Randomized action sequences $\mathbf{a}_{t:t+T}$ within the robot's FOV are generated and executed until a collision or a timeout occurs. Datapoints consisting of $d = (\mathbf{x}_t, \mathbf{s}_t, \mathbf{a}_{t:t+T}, \mathbf{\hat{c}}^{col}_{t+1:t+T+1})$ are collected. The data is also augmented by performing a horizontal flip and appending random actions at the end of the collision episodes, similar to~\cite{ORACLE_ICRA2022}. The depth images $\mathbf{x}_t$ from this are passed through the learned \ac{sevae} model to obtain its latent representation $\bm{\mu}_t$. Finally, datapoints  $d^\prime = (\bm{\mu}_t, \mathbf{s}_t, \mathbf{a}_{t:t+T}, \mathbf{\hat{c}}^{col}_{t+1:t+T+1})$ are created replacing the depth image $\mathbf{x}_t$ with its corresponding latent representation $\bm{\mu}_t$ to train the \ac{cpn}.



%%%%%%%%%%%%%%%%%%%%%%%%%%%%%%%%%%%%%%%%%%%%%%%%%%%%%%%%%%%%%%%%%
\begin{figure}[h]
\includegraphics[width=\columnwidth]{figures/ReconstructionComparisonv2_seVAE.png}
\vspace{-4ex}
\caption{The proposed method is compared with Fast Fourier Transform (FFT) and vanilla-VAE without semantic weighted loss on real (1) and simulated images (2,3). The FFT reconstruction is created from the largest $64$ magnitudes in the frequency domain representation, while the vanilla-VAE and the semantically-enhanced VAE use a $128$ dimensional latent vector. Ellipsoids are drawn around expected reconstructions of thin obstacles from the input depth images.}
\label{fig:vae_fft_reconstruction_comparison}
\vspace{-4ex}
\end{figure}
%%%%%%%%%%%%%%%%%%%%%%%%%%%%%%%%%%%%%%%%%%%%%%%%%%%%%%%%%%%%%%%%%



\subsection{Comparison of reconstruction methods for thin obstacles}\label{sec:reconstruction_comparison}
First, the proposed semantically-enhanced VAE is evaluated by comparing it against a baseline compression based on the Fast Fourier Transform (FFT) and a vanilla-VAE trained without the weighing function $\lambda(\mathbf{x}_{seg})$. To allow for a fair comparison with the $128$ dimensional latent variable $\bm{\mu}$ of the VAEs, the FFT reconstructions are generated from the $64$ complex frequency-domain representations having the largest magnitudes while zeroing out the others. Reconstructions from the compared methods are depicted in Figure~\ref{fig:vae_fft_reconstruction_comparison}. 
The FFT reconstruction is unable to preserve information at this resolution especially for the features corresponding to smaller-sized obstacles. Additionally, the presence of sensor noise in real images degrades the performance of the FFT even further. The vanilla-VAE is able to reconstruct larger parts of the images well but misses out on the regions with small cross-section (thin) obstacles for both real and simulated images. The \ac{sevae} is able to reconstruct regions with smaller cross sections better than the above approaches. The key role of semantically-augmented training is particularly visible in Figure~\ref{fig:vae_fft_reconstruction_comparison}. We statistically compare the performance of the three approaches to consider the MSE over the whole image and also specifically over only the semantically labeled pixels (thin obstacles) of an image and present the results in Table~\ref{tab:mse_comparison_table_fft_vanilla_vae}. For comparison, we normalize the valid pixel values between $0$ and $1$. As shown, \ac{sevae} presents a significantly better performance in reconstructing the semantic regions, with a relatively small reduction in the performance over the whole image against the vanilla-VAE using the same number of parameters, while also outperforming FFT-based compression. It is noted that relatively inferior reconstruction quality is less significant for large objects, as long as collision-avoidance is concerned, but it is critical not to miss thin obstacles. 


\begin{table}[h]
    \vspace{-2ex}
    \centering
    \caption{Comparison of MSE for reconstructed images with different methods (seVAE, vanilla-VAE, FFT-based).}
    \vspace{-2ex}
    \begin{tabularx}{\columnwidth}{|l|X|X|X|}
      \hline
      \multicolumn{4}{|c|}{\textbf{Simulated Images (Count: $2239$)}} \\
      \hline
      \textbf{MSE over:} & \textbf{FFT} & \textbf{vanilla-VAE} & \textbf{seVAE}\\
       \hline
      Entire image & $553.21$ & $\mathbf{276.64}$ & $404.50$  \\
      \hline
      Semantic pixels  & $82.14$  & $89.51$  & $\mathbf{22.13}$  \\
      \hline
      \multicolumn{4}{|c|}{\textbf{Real Images (Count: $365$)}} \\
      \hline
      \textbf{MSE over:} & \textbf{FFT} & \textbf{vanilla-VAE} & \textbf{seVAE}\\
      \hline
      Entire image & $2804.28$ & $\mathbf{350.01}$ & $420.35$  \\
      \hline
      Semantic pixels  & $156.95$  & $69.16$  & $\mathbf{28.41}$  \\
      \hline
    \end{tabularx}
    \label{tab:mse_comparison_table_fft_vanilla_vae}
    \vspace{-4ex}
\end{table}


%%%%%%%%%%%%%%%%%%%%%%%%%%%%%%%%%%%%%%%%%%%%%%%%%%%%%%%%%%%%%%%%% 
\begin{figure}[h!]
\centering
\includegraphics[width=0.9\columnwidth]{figures/vae_sim_v6.png}
\caption{Results from one of the simulation missions conducted in a Gazebo environment with diverse obstacle configurations including thin obstacles (cross section $4~\textrm{cm}$). The sub-figures 1 and 2 show two instances during the mission detailing a) the point cloud, b) the onboard camera image both overlaid with markers representing estimated endpoints of collision-free action sequences, c) the raw depth image, and d) reconstructed depth image from \ac{sevae}. Sub-figure 3 shows the final map and the trajectory followed by the robot. The simulation shows the ability of the \ac{sevae} to reconstruct thin and regular objects well, and the overall method to avoid obstacles.}
\vspace{-4ex}
\label{fig:vae_sim}
\end{figure}
%%%%%%%%%%%%%%%%%%%%%%%%%%%%%%%%%%%%%%%%%%%%%%%%%%%%%%%%%%%%%%%%%

\subsection{Simulation Studies}\label{sec:sim_studies}
Subsequently, the proposed method is evaluated in complex simulation studies involving diverse obstacle configurations including distributed thin objects such as rods. We design three simulation worlds consisting of three sections each, with different kinds of obstacles placed using Poisson disc sampling. The first section contains exclusively large-sized obstacles with Poisson disc sampling radius $r_1$, the second section contains both large sized, and thin obstacles with cross section $4~\textrm{cm}$, sampled independently of each other with radii $r_2$ and $r_3$ respectively, and the third section consists only of thin rods with a similar cross section sampled with radius $r_4$. Each section is $50~\textrm{m}\times50~\textrm{m}$ and placed serially. The robot is commanded to travel $150~\textrm{m}$ along the course at a speed of $1.0~\textrm{m/s}$. To derive the UT sigma points, we use $\boldsymbol{\Sigma}_t = \textrm{diag}({\sigma_v}^2, {\sigma_v}^2, {\sigma_v}^2, 0, 0, 0)$, with ${\sigma_v}=0.2~\textrm{m}/\textrm{s}$. We perform $20$ runs of the experiment per environment for both the proposed method (trained with datapoints $d^\prime$) and a $3$D extension of our end-to-end ORACLE method~\cite{ORACLE_ICRA2022} (trained with datapoints $d$) with different initial positions and orientations. A few instances and the full path from one of the missions of the proposed method are shown in Figure~\ref{fig:vae_sim} and the statistical results are logged in Table~\ref{tab:sim_result}. The environment sampling variables are listed beside each environment in the format ($r_1$, $r_2$, $r_3$, $r_4$). It is noted that the ensemble of neural networks is not used for ORACLE to have a fair comparison. 
The proposed method outperforms ORACLE with a higher success rate in terms of completing the course without collision, in denser environments.



\begin{table}[h]
 \vspace{-2ex}
 \caption{Comparison between success rates of end-to-end vs proposed method in simulation.}
 \vspace{-2ex}
\centering
\begin{tabular}{|c|c|c|}  
\hline
\textbf{Environment ($r_1$, $r_2$, $r_3$, $r_4$)} & \textbf{Method} & \textbf{Success \%} \\
\hline
    \multirow{2}{*}{Sparse ($6.5$, $6.5$, $4.5$, $3.5$)}
       & Proposed method & $95\%$  \\
        & End-to-end ORACLE        & $95\%$     \\
\hline
     \multirow{2}{*}{Medium ($6.25$, $6.25$, $3.5$, $3.0$)} & Proposed method          & $\mathbf{80}\%$ \\
      & End-to-end ORACLE     & $75\%$     \\            
\hline
     \multirow{2}{*}{Dense ($6.0$, $6.0$, $3.0$, $2.5$)} & Proposed method          & $\mathbf{60}\%$   \\
       & End-to-end ORACLE     & $45\%$  \\
\hline
\end{tabular}
\label{tab:sim_result}
 \vspace{-4ex}
\end{table}

\begin{figure*}[ht!]
\includegraphics[width=\textwidth]{figures/vae_experiment_v12.png}
\vspace{-4ex}
\caption{Results from the two real-world experiments conducted to evaluate the proposed method. The first took place in a long cluttered corridor ($\sim 50~\textrm{m}$) including thin obstacles as small as $4.8~\textrm{cm}$ in width. The second experiment was conducted in a relatively smaller corridor ($\sim 20~\textrm{m}$) with predominantly thin obstacles, some as small as $2.5~\textrm{cm}$ in width, some with texture-less reflective surfaces. The sub-figures 1.1 and 1.2 show two instances during experiment 1 (similarly 2.1, 2.2 for experiment 2) detailing a) $3^{\textrm{rd}}$ person view of the robot and the environment, b) onboard color camera image overlaid with markers representing estimated endpoints of collision-free action sequences, c) raw depth image, and d) reconstructed depth image from \ac{sevae}. 1.3 (similarly 2.3 for experiment 2) shows the final map and trajectory followed by the robot in experiment 1. The reconstructed maps are only for visualization purposes and are not computed onboard the robot during flights. It can be seen that the method is successfully able to respond to and avoid all the obstacles.}
\vspace{-2ex}
\label{fig:experiment_diagram}
\end{figure*}
%%%%%%%%%%%%%%%%%%%%%%%%%%%%%%%%%%%%%%%%%%%%%%%%%%%%%%%%%%%%%%%%%

\subsection{Experimental Evaluation}\label{sec:experimental_evaluation}

Beyond simulations, the proposed method was evaluated in two real-world confined environments further involving thin obstacles using a collision-tolerant aerial robot design similar to \cite{rmfowl_icuas} which integrates a sensor suite including an Intel Realsense D455 RGB-D sensor, an Intel Realsense T265 visual-inertial module, an mRo Pixracer flight controller, and an NVIDIA Xavier NX board. To derive the UT sigma points, we use $\boldsymbol{\Sigma}_t = \textrm{diag}({\sigma_v}^2, {\sigma_v}^2, {\sigma_v}^2, 0, 0, 0)$, with ${\sigma_v}=0.2~\textrm{m}/\textrm{s}$. One experiment is performed (with an average robot speed of $1.0~\textrm{m/s}$) in a long cluttered corridor with thin obstacles (as small as $4.8~\textrm{cm}$ in width) obstructing the path of the robot. The goal point is defined $50~\textrm{m}$ from the starting point along the corridor and expressed as a goal direction for the navigation method exploiting only a partial state without position information. A second experiment is conducted (with an average robot speed of $0.75~\textrm{m/s}$) in a smaller environment with challenging hard-to-perceive obstacles with small cross-sections (as small as $2.5~\textrm{cm}$ in width), and texture-less reflective surfaces. The size and the material of the obstacles in these experiments pose challenges for accurate depth reconstruction using the depth camera onboard the robot. The resolution of the RealSense D455 is set to $480\times640~\textrm{pixels}$ and the depth image is downsampled to the input resolution of the encoder ($270\times480~\textrm{pixels}$) before encoding. The inference times of \ac{sevae} and \ac{cpn} on the platform are $11~\textrm{ms}$ and $29~\textrm{ms}$ respectively. The results of these experiments are shown in Figure~\ref{fig:experiment_diagram}.

%%%%%%%%%%%%%%%%%%%%%%%%%%%%%%%%%%%%%%%%%%%%%%%%%%%%%%%%%%%%%%%%% 
\begin{figure}[ht!]
\centering
\includegraphics[width=0.9\columnwidth]{figures/VAE_vs_end_to_end_comparisons_v7_corrected.png}
\caption{Comparison of collision predictions of ORACLE end-to-end trained purely in simulation against the proposed method with un-filtered depth data in the real-world experiments. Our proposed method predicts collision-free action sequences and selects the safest one toward the goal whereas ORACLE incorrectly predicts the collision-free action sequences.}
\vspace{-3ex}
\label{fig:vae_vs_end_to_end_oracle}
\end{figure}
%%%%%%%%%%%%%%%%%%%%%%%%%%%%%%%%%%%%%%%%%%%%%%%%%%%%%%%%%%%%%%%%%


Finally, the proposed approach is also experimentally compared against the end-to-end trained ORACLE as in Section~\ref{sec:sim_studies}. Specifically, the collision prediction step of ORACLE was ran for the depth frames of the above experiments. Notably, originally ORACLE also runs a filter that attempts to fill-in invalid pixels coming from the sensor~\cite{ku2018defense}, which is however computationally burdensome ($\sim20~\textrm{ms}$ compute time on an NVIDIA Xavier NX) and not needed in the proposed method thus not executed for the purposes of fair comparison. As shown in Figure~\ref{fig:vae_vs_end_to_end_oracle}, the end-to-end ORACLE incorrectly predicts action sequences --that may cause the robot to collide with the environment-- as collision-free, unlike the proposed approach which enables safe action sequences towards the goal direction. Despite the potential benefits of end-to-end training, ORACLE is unable to incorporate the perceptual focus on thin obstacles of the new exploit using the \ac{sevae} and is also unable to exploit real sensor observations in the training. As a result, learning to avoid hard-to-perceive thin obstacles is only trained from simulated data which present significant differences compared to real data (not merely white noise and other easy-to-simulate errors)~\cite{yang2018segstereo,ahn2019analysis,mallick2014characterizations} thus widening the sim-to-real gap. Note that, the ensemble of neural networks in~\cite{ORACLE_ICRA2022} is not used to have a fair comparison.  
